\AliusParagraph{The video and audio version of this episode are available on}

\AliusParagraph{Conscious experiences during sleep have long been considered mysterious and typically difficult to assess as they can only be shared after awakening. Even the memories that we are left with in the morning are often faint and sometimes bizarre. So a question naturally arises : do memories after awakening truly reflect the type of conscious experiences that we have during sleep ? Even more dramatically, could memories be completely made up to such an extent that we actually do not dream at all during sleep? (Malcolm, 1953; Dennett, 1976). This hypothesis has been debated over the past 50 years but the scientific consensus is now suggesting that dream reports are actually faithful (Windt, 2013). One of the strongest arguments in favor of this hypothesis comes from people who have a pathological condition called REM behavior disorder. Because people with REM behavior disorder are unable to inhibit their movements while dreaming, researchers could observe their actions and found that they fit remarkably well with the dream experiences that were described after awakening (Oudiette et al., 2009).}

\AliusParagraph{Could memories be made up to the such an extent}

\AliusPullQuote{that we actually do not dream at all during sleep?}

\AliusParagraph{So if we can trust dream reports, what do they tell us? Researchers have studied this question systematically by waking up participants at different moments of the night and asking them to describe the conscious experience that they recall from their dreams. What they found is that dream contents varied greatly depending on the time of the night (Casagrande et al., 1996; Siclari et al., 2013). While we are falling asleep, we can experience intense conscious activity that is called hypnagogic dreams. These dreams often reflect what we have been experiencing during the day. This is known as the Tetris effect, because doing some actions for several hours, such as playing the video game Tetris, causes images of this game to be replayed during sleep (Stickgold, 2000). Even while firmly asleep, our dreams tend to be composed of images and actions related to our daily activities. But as the night goes on, the more complex our dreams become.}

\AliusParagraph{One aspect of dreams that has drawn particular attention is our experience of the self (Revonsuo, 2005). Our subjective experience can be in the first person perspective, embracing the point of view of our body. However, we may also experience dreams in the third person, floating slightly above ourselves. We may even be absent from the dream, adopting a disembodied point of view that simply watches the dream unfold. Another focus on the dream experience has been the perception of time. While a dream report can describe actions taking place over several hours or days, we only dream for a few minutes or hours. However, a good correlation was found between the actual time spent asleep and the duration of dream experiences (Dement \& Wolpert, 1958). An explanation for this fact could be that, in our dreams, we are not so surprised by transitions that involve jump in time, a phenomenon that is also found in books or movies and called narrative ellipses.}

\AliusPullQuote{The dreamless sleep experience has been\textbackslash{}\textbackslash{}described in particular detail by great meditators.}

\AliusParagraph{Finally, we can also wake up knowing that we experienced something in our dream but without being able to say precisely what it was about. While this may reflect a failure of our memory, this could also tap into the existence of conscious experiences that are lacking any content. This phenomenon is known as mind-blanking during wakefulness (Ward \& Wegner, 2012) and white dreams during sleep (Fazekas et al., 2019). This dreamless sleep experience has been described in particular detail by great meditators who have trained themselves to remain conscious even after falling asleep (Windt et al., 2016).}

\AliusParagraph{Another trainable skill that has sparkled the interest of researchers as a tool to probe the nature of dreams is lucid dreaming. Lucid dreaming is the ability to become aware of dreaming while in a dream (LaBerge \& Rheingold, 1990). A question that has been investigated is to what extent the body of the sleeper is involved when we dream. Researchers found that sleepers' respiratory rhythm slows down if lucid dreamers are asked to stop breathing in their dream (Oudiette et al., 2019). By asking different questions to lucid dreamers while recording their body and brain responses, five laboratories around the world have now succeeded in communicating in real-time with lucid dreamers (Konkoly et al., 2021). This rise in research about dream and lucid dreaming offers us novel and unique perspectives to better understand what it is like to be conscious while sleeping.}

\AliusPullQuote{Lucid dreaming is the ability to become aware\textbackslash{}\textbackslash{}to become aware of dreaming while in a dream.}

\AliusSectionHeading{References}

\AliusReferenceEntry{Casagrande, M., Violani, C., Lucidi, F., Buttinelli, E., \& Bertini, M. (1996).}

\AliusReferenceEntry{Variations in Sleep Mentation as a Function of Time of Night.}

\AliusReferenceEntry{International Journal of Neuroscience, 85(1–2), 19–30.}

\AliusReferenceEntry{Dement, W., \& Wolpert, E. A. (1958). The relation of eye movements, body}

\AliusReferenceEntry{motility, and external stimuli to dream content. Journal of Experimental}

\AliusReferenceEntry{Psychology, 55(6), 543–553.}

\AliusReferenceEntry{Dennett, D. C. (1976). Are dreams experiences?. The Philosophical Review, 85(2),}

\AliusReferenceEntry{151-171.}

\AliusReferenceEntry{Fazekas, P., Nemeth, G., \& Overgaard, M. (2019). White dreams are made of}

\AliusReferenceEntry{colours: What studying contentless dreams can teach about the neural}

\AliusReferenceEntry{basis of dreaming and conscious experiences. Sleep Medicine Reviews, 43,}

\AliusReferenceEntry{84–91.}

\AliusReferenceEntry{Konkoly, K. R., Appel, K., Chabani, E., Mangiaruga, A., Gott, J., Mallett, R.,}

\AliusReferenceEntry{Caughran, B., Witkowski, S., Whitmore, N. W., Mazurek, C. Y., Berent, J.}

\AliusReferenceEntry{B., Weber, F. D., Türker, B., Leu-Semenescu, S., Maranci, J.-B., Pipa, G.,}

\AliusReferenceEntry{Arnulf, I., Oudiette, D., Dresler, M., \& Paller, K. A. (2021). Real-time}

\AliusReferenceEntry{dialogue between experimenters and dreamers during REM sleep. Current}

\AliusReferenceEntry{Biology, S0960982221000592.}

\AliusReferenceEntry{LaBerge, S., \& Rheingold, H. (1990). Exploring the World of Lucid Dreaming. New}

\AliusReferenceEntry{York: Ballantine Books.}

\AliusReferenceEntry{Malcolm, N. (2017). Dreaming. Routledge.}

\AliusReferenceEntry{Oudiette, D., De Cock, V. C., Lavault, S., Leu, S., Vidailhet, M., \& Arnulf, I.}

\AliusReferenceEntry{(2009). Nonviolent elaborate behaviors may also occur in REM sleep}

\AliusReferenceEntry{behavior disorder. Neurology, 72(6), 551–557.}

\AliusReferenceEntry{Revonsuo, A. (2005). The self in dreams. In The lost self: Pathologies of the brain and}

\AliusReferenceEntry{identity (pp. 206–219). Oxford University Press Oxford.}

\AliusReferenceEntry{Siclari, F., LaRocque, J. J., Postle, B. R., \& Tononi, G. (2013). Assessing sleep}

\AliusReferenceEntry{consciousness within subjects using a serial awakening paradigm. Frontiers}

\AliusReferenceEntry{in Psychology, 4.}

\AliusReferenceEntry{Stickgold, R. (2000). Replaying the Game: Hypnagogic Images in Normals and}

\AliusReferenceEntry{Amnesics. Science, 290(5490), 350–353.}

\AliusReferenceEntry{Windt, J. M. (2013). Reporting dream experience: Why (not) to be skeptical about}

\AliusReferenceEntry{dream reports. Frontiers in Human Neuroscience, 7.}

\AliusReferenceEntry{Windt, J. M., Nielsen, T., \& Thompson, E. (2016). Does Consciousness Disappear}

\AliusReferenceEntry{in Dreamless Sleep?Trends in Cognitive Sciences, 20(12), 871–882.}

\AliusReferenceEntry{Ward, A. F., \& Wegner, D. M. (2013). Mind-blanking: When the mind goes away.}

\AliusReferenceEntry{Frontiers in psychology, 4, 650.}

\AliusReferenceEntry{Oudiette, D., Dodet, P., Ledard, N., Artru, E., Rachidi, I., Similowski, T., \&}

\AliusReferenceEntry{Arnulf, I. (2018). REM sleep respiratory behaviours match mental content}

\AliusReferenceEntry{in narcoleptic lucid dreamers. Scientific reports, 8(1), 1-10.}
